\documentclass[10pt,a4paper]{article}
\usepackage[T1]{fontenc}
\usepackage[utf8]{inputenc}
\usepackage{lmodern}
\usepackage{amsmath,amssymb}
\usepackage[margin=21mm]{geometry}
\usepackage{microtype}
\usepackage[hidelinks]{hyperref}
\usepackage{enumitem}
\setlist[itemize]{leftmargin=*,nosep,topsep=3pt}
\setlength{\parindent}{0pt}
\setlength{\parskip}{5pt}
\setlength{\emergencystretch}{3em}
\newcommand{\question}[1]{\par\medskip{\normalsize\bfseries #1}\par\smallskip}
\pagestyle{plain}

\begin{document}
{\Large\bfseries Dopamine and deep brain stimulation accelerate the neural dynamics of volitional action in Parkinson's disease\par}

\textbf{Paper:} Koehler et al. (2024).\quad Published human electrophysiology study.\par
\href{https://doi.org/10.1093/brain/awae219}{DOI: \nolinkurl{10.1093/brain/awae219}}\par
{\footnotesize Page references below use PDF page numbers, including front matter.}

\question{To which general question does this paper relate?}
Measurement, impaired movement initiation, and the control objective. Relates treatment effects to the interval between preparatory neural activity and muscle activation.

\question{How does it answer the general questions?}
\textbf{Measured.} Sensorimotor ECoG, STN-LFP and EMG were recorded during self-initiated movements in 25 patients. Readiness potentials and classifiers estimated preparatory activity; directional spectral analyses assessed cortico-subthalamic communication (PDF pp.~1--8).

\textbf{Treatment effect.} Dopaminergic medication and STN-DBS shortened preparation-to-execution latencies and shifted directional coupling towards theta. These findings associate treatment benefit with changes in the timing and frequency structure of communication; they do not establish beta-power minimization as a sufficient objective (PDF pp.~8--11).

\textbf{Model relevance.} Movement-initiation latency is a useful behavioural outcome. These latencies should not be substituted directly for axonal delays in the notebook model: they summarize distributed preparation and execution processes. No nonlinear delayed-state model or feedback controller is identified.

\textbf{Version.} This is the published counterpart of note~05; the two files are not independent studies.

\question{What papers does it cite that could provide more information?}
\begin{itemize}
\item Kornhuber and Deecke (1965), work on the \emph{Bereitschaftspotential}: origin of the preparatory-potential measure (ref.~2, PDF p.~11).
\item Lofredi et al. (2018), \emph{Dopamine-dependent scaling of subthalamic gamma bursts with movement velocity in patients with Parkinson's disease}: another movement-related outcome (ref.~56, PDF p.~12).
\end{itemize}

\end{document}
